\documentclass[14pt]{beamer}
\usetheme{defense}
\graphicspath{{../../figures/}}


\usepackage{tikz}
\usepackage{svg}
\usepackage{anyfontsize}
\setcounter{tocdepth}{1}

\title{Memristors-based recurent modules for neural computing}

\subtitle{}

\newcommand {\Supervisor} {
  {Dr. Diogo Caetano} and {Dr. Ruxandra Barbulescu}}
\newcommand {\CommitteeMembers} {
  {Prof./Dr. Lorem Ipsum}}
\newcommand {\Chairperson} {{Prof./Dr. Lorem Ipsum}}
\author[V. BARBAZA]{{Valentin BARBAZA}}

\date{17 november 2023}

\logo{
  \begin{tikzpicture}[overlay,remember picture]
    \node[left=0cm] at (current page.33){
      \includegraphics[height=1cm]{logos/inesc-mn.png}
      \hspace{1pt}
      \includegraphics[height=1cm]{logos/ist.eps}
      \includegraphics[height=1cm]{logos/inesc-id.eps}
    };
  \end{tikzpicture}
}

\begin{document}

\frame{\titlepage}

\begin{frame}
  \frametitle{Table of Contents}
  \tableofcontents
\end{frame}

\section{Motivation and objectives}

\begin{frame}{\insertsection}
  \begin{columns}
    \column{.6\textwidth}
    \begin{itemize}
      \item NNs complexity grows rapidly
      \item So does its computation time
    \end{itemize}
    \column{.4\textwidth}
    \includesvg[width=\columnwidth,pretex=\scriptsize]{NN_explained.svg}
  \end{columns}
\end{frame}

\begin{frame}{\insertsection}
  \begin{columns}
    \column{.6\textwidth}
    \begin{itemize}
      \item Reduce computation by using analog
      \item Allows use for time sensitive problems
    \end{itemize}
    \column{.4\textwidth}
    \includesvg[width=\columnwidth,pretex=\scriptsize]{NN_explained.svg}
  \end{columns}
\end{frame}

\section{Introduction}

\begin{frame}{\insertsection}{RNNs}
  \begin{columns}
    \column{.4\textwidth}
    \includesvg[width=\columnwidth,pretex=\small]{rnn/rnnFolded}
    \column{.6\textwidth}
    RNNs are a type of NN :
    \begin{itemize}
      \item Used for time dependent data
      \item That have a feedback connection
    \end{itemize}
  \end{columns}
  \vspace{0.8cm}

  \begin{equation}\label{eq:rnn}
    \overrightarrow{h_t}=f(\overrightarrow{x_t},\overrightarrow{h_{t-1}})
  \end{equation}
\end{frame}

\begin{frame}{\insertsection}{LSTMs}
  \includesvg[width=\columnwidth,pretex=\scriptsize]{lstm/lstmCell}
\end{frame}

\begin{frame}{\insertsection}{GRUs}
  \includesvg[width=\columnwidth,pretex=\scriptsize]{gru/encoderCell}
\end{frame}

\begin{frame}{\insertsection}{Memristors}
  \begin{columns}
    \column{.65\textwidth}
    \includesvg[width=\columnwidth,pretex=\tiny]{memristor/memristor}
    \column{.4\textwidth}
    \begin{itemize}
      \item Recent electical component
      \item Variable resistance
    \end{itemize}
  \end{columns}
\end{frame}


\section{The circuits}
\subsection{Activation functions}
\begin{frame}{\insertsection}{\insertsubsection}
  \begin{columns}
    \column{.7\textwidth}
    \includesvg[width=\columnwidth,pretex=\scriptsize]{activation/afCircuit}
    \column{.4\textwidth}
    \begin{itemize}
      \item Creates analog activation functions
      \item Is used for both tanh and sigmoid activation functions
    \end{itemize}
  \end{columns}
\end{frame}

\begin{frame}{\insertsection}{\insertsubsection}
  \begin{center}
    \includesvg[width=0.75\columnwidth,pretex=\tiny]{activation/afGraph}
  \end{center}
\end{frame}

\subsection{Memory cell}

\begin{frame}{\insertsection}{\insertsubsection}
  \begin{columns}
    \column{.65\textwidth}
    \includesvg[width=\columnwidth,pretex=\tiny]{memcell/memCircuit}
    \column{.4\textwidth}
    \begin{itemize}
      \item Stores analog value
      \item Has two CMOS switches to avoid leakage
    \end{itemize}
  \end{columns}
\end{frame}

\begin{frame}{\insertsection}{\insertsubsection}
  \begin{center}
    \includesvg[width=.75\columnwidth,pretex=\tiny]{memcell/data-loss}
  \end{center}
\end{frame}

\subsection{Crossbar circuit}

\begin{frame}{\insertsection}{\insertsubsection}
  \begin{columns}
    \column{.4\textwidth}
    \begin{itemize}
      \item Performs analog VMM
      \item Can be parallel or serialized
    \end{itemize}
    \column{.65\textwidth}
    \includesvg[width=\columnwidth, pretex=\tiny]{crossbar/crossbarUse}
  \end{columns}
\end{frame}

\begin{frame}{\insertsection}{\insertsubsection}
  \includesvg[width=\columnwidth, pretex=\tiny]{crossbar/doubleMem}
  \begin{equation}
    \label{eq:doubleMem2}
    y_0=R_f\cdot\sum_{k=0}^n(\sigma_{k+}-\sigma_{k-})\cdot x_k
  \end{equation}
\end{frame}

\subsection{LSTM circuit}

\begin{frame}{\insertsection}{\insertsubsection}
  \begin{center}
    \includesvg[width=1\columnwidth,pretex=\fontsize{4}{5}\selectfont]{lstm/lstmCircuit}
  \end{center}
\end{frame}


\section{The tools}
\subsection{Weight generation}
\begin{frame}{\insertsection}{\insertsubsection}
  The weights are trained :
  \hfill\includegraphics[width = 4cm]{tensorflowKeras.jpg}
  \begin{itemize}
    \item Using Keras framework implemented in tensorflow python module
    \item In software and exported to a file
    \item Using the analog activation functions
  \end{itemize}
\end{frame}

\subsection{Netlist generation}
\begin{frame}{\insertsection}{\insertsubsection}
  I created a script capable of generating modular netlist. It can take in :
  \begin{itemize}
    \item Any number of time steps
    \item Any number of inputs
    \item Any serial size
    \item Dense layers
    \item LSTM layers
    \item GRU layers
  \end{itemize}
\end{frame}

\section{The datasets}
\subsection{Airline}
\begin{frame}{\insertsection}{\insertsubsection}
  The airline dataset contains the monthly number of passengers on international airlines from 1949 to 1960.
  \includesvg[width=\columnwidth,pretex=\tiny]{datasets/airline}
\end{frame}

\subsection{C. elegans}
\begin{frame}{\insertsection}{\insertsubsection}
  \begin{columns}
    \column{.35\textwidth}
    This dataset is more specific :
    \begin{itemize}
      \item It has has 4 inputs
      \item Across 1000 time steps
      \item Contains 40 sequences
    \end{itemize}
    \column{.7\textwidth}
    \includesvg[width=.95\columnwidth,pretex=\tiny]{datasets/celegans/ioValid0}
  \end{columns}
\end{frame}

\section{Results}
\subsection{Airline}
\begin{frame}{\insertsection}{\insertsubsection}
  \begin{columns}
    \column{.68\textwidth}
    \includesvg[width=\columnwidth,pretex=\tiny]{results/airline/analog}
    \column{.37\textwidth}
    \begin{itemize}
      \item Better results for $n_s=1$
      \item Very promising looking
    \end{itemize}
  \end{columns}
\end{frame}

\subsection{C. elegans}

\begin{frame}{\insertsection}{\insertsubsection}
  \begin{center}
    \includesvg[width=0.8\columnwidth,pretex=\tiny]{results/celegans/out5}
  \end{center}
\end{frame}

\begin{frame}{\insertsection}{\insertsubsection}
  \begin{center}
    \includesvg[width=\columnwidth,pretex=\tiny]{results/celegans/smooth5}
  \end{center}
\end{frame}

\section{Conclusion}

\begin{frame}{\insertsection}
  \begin{itemize}
    \item Decent results
    \item Very quick results
    \item Instability when dealing with large amount of time steps
  \end{itemize}
\end{frame}

\end{document}
